\clearpage
\fancyhead[L]{\small GAUSS AN BODE.}
\SourceLeaf{454}

\begin{center}
{\large GAUSS AN BODE.}
\end{center}

\begin{center}
\rule{6cm}{0.4pt}\\
Astronomisches Jahrbuch für 1827. Seite 128..130. Berlin 1824.\\
\rule{6cm}{0.4pt}
\end{center}

\begin{flushright}
Göttingen, den 21. April 1824.
\end{flushright}

Bei all diesen Beobachtungen war der Komet so lichtschwach, dass er gar keine Fadenbeleuchtung vertrug. Komet und Fadennetz ließen sich daher nur sichtbar machen, indem man das Licht abwechselnd abschirmte und wieder freigab. Unter diesen Umständen war bei den Beobachtungen, besonders bei der Bestimmung der Rektaszensionen, nur eine geringe Genauigkeit möglich. Die angegebenen Zeiten beziehen sich auf das Mittel der beobachteten oder vielmehr geschätzten Fadenantritte. An Tagen, an denen nicht an allen Fäden beobachtet wurde, weichen sie daher etwas von der wahren Kulminationszeit ab.

\begin{center}
\begingroup\small\setlength{\tabcolsep}{4pt}
\begin{tabular}{lr|c|c|c}
\multicolumn{2}{c|}{1824} & \shortstack{Mittlere Zeit\\in Göttingen} & Rektaszension & Deklination \\
\hline
Febr. & 31 & $13^{\mathrm u}\,28^{\mathrm m}\,20^{\mathrm s}$ & $152^\circ\,16'\,29''$ & $72^\circ\,38'\,35''$ N. \\
Jan. & 3 & $11\quad55\quad7$ & $131\quad23\quad10$ & $69\quad14\quad9$ \\
 & 7 & $10\quad38\quad3$ & $116\quad12\quad50$ & $63\quad14\quad56$ \\
 & 8 & $10\quad24\quad5$ & $113\quad50\quad49$ & $61\quad46\quad15$ \\
 & 19 & $8\quad52\quad23$ & $101\quad49\quad31$ & $48\quad48\quad51$ \\
 & 20 & $8\quad48$ & ---\quad---\quad--- & $47\quad55$ \quad : \\
 & 27 & $8\quad11\quad46$ & $99\quad26\quad51$ & $42\quad48\quad49$ \\
 & 28 & $8\quad6\quad47$ & $99\quad18\quad8$ & $42\quad12\quad46$ \\
März & 2 & $7\quad53\quad59$ & $99\quad0\quad39$ & $40\quad32\quad3$ \\
\end{tabular}
\endgroup
\end{center}

Einige frühere Beobachtungen, die ich außerhalb des Meridians mit einem Kreismikrometer oder Heliometer gemacht habe, habe ich noch nicht reduziert. Ich muss nämlich die Vergleichssterne, die in der Histoire Céleste fehlen, erst selbst bestimmen.

Herr Dr. \textsc{Schmidt}, ein junger Mann mit ausgezeichneten mathematischen Kenntnissen, hat aus diesen und einigen anderen Beobachtungen die folgenden parabolischen Elemente berechnet:

\begin{center}
\begingroup\small
\begin{tabular}{p{.51\linewidth}p{.41\linewidth}}
Länge des Knotens \mbox{.\,.\,.\,.\,.\,.\,.\,.\,.\,.\,.\,.} & $303^\circ\,3'\,51''3$ \\
Neigung der Bahn \mbox{.\,.\,.\,.\,.\,.\,.\,.\,.\,.\,.\,.} & $76\quad12\quad14.0$ \\
Länge der Sonnennähe \mbox{.\,.\,.\,.\,.\,.\,.\,.\,.\,.\,.} & $274\quad34\quad14.4$ \\
Durchgang durch dieselbe \mbox{.\,.\,.\,.} & 1823 Dezember 9.47459 Göttinger Zeit. \\
Logarithmus des kleinsten Abstands. \mbox{.\,.\,.\,.\,.\,.} & 9.3553041 \\
\multicolumn{2}{l}{Rückläufige Bewegung.} \\
\end{tabular}
\endgroup
\end{center}

\Needspace{23\baselineskip}
Der Vergleich dieser Elemente mit den bisher bekannt gewordenen Beobachtungen ergab die folgenden Unterschiede:

\Needspace{12\baselineskip}
\SourceLeaf{455}


\begin{center}
\begin{tabular}{lr|r|r|l}
 & & \shortstack{in\\Rektaszension} & \shortstack{in\\Deklination} & Beobachter \\
\hline
Jan. & 1 & 0 & 0 & Bouvard \\
 & 3 & $+12.5$ & $-22.6$ & Nicolai \\
 & 4 & $+13.6$ & $+1.6$ & Nicolai \\
 & 5 & $+34.5$ & $-6.5$ & Nicolai \\
 & 11 & $+0.2$ & $+9.5$ & Olbers \\
 & 14 & $+27.5$ & $+0.3$ & Olbers \\
 & 24 & $+48.6$ & $+3.4$ & Schumacher \\
 & 25 & $+85.0$ & $-3.3$ & Schumacher \\
 & 26 & $+44.0$ & $+11.0$ & Soldner \\
 & 27 & $+33.0$ & $+10.4$ & Soldner \\
 & 30 & $-26.5$ & $+20.7$ & Schumacher \\
 & 31 & $-36.7$ & $+25.5$ & Soldner \\
 & 31 & $+24.9$ & $+14.5$ & Gauss \\
Febr. & 2 & $+963.4$ & $-5.7$ & Schumacher \\
 & 3 & $-28.9$ & $+12.0$ & Schumacher \\
 & 3 & $-40.6$ & $-8.6$ & Gauss \\
\end{tabular}
\end{center}

\begin{center}
\begin{tabular}{lr|r|r|l}
 & & \shortstack{in\\Rektaszension} & \shortstack{in\\Deklination} & Beobachter \\
\hline
Febr. & 7 & $-4.5$ & $+18.5$ & Gauss \\
 & 8 & $-40.6$ & $-2.0$ & Gauss \\
 & 11 & $-18.4$ & $-22.0$ & Schumacher \\
 & 13 & $-10.4$ & $-1.2$ & Schumacher \\
 & 18 & $-2.6$ & $-66.1$ & Olbers \\
 & 19 & $-1.0$ & $+1.0$ & Gauss \\
 & 19 & $+26.0$ & $-26.7$ & Olbers \\
 & 21 & $+36.2$ & $-38.8$ & Olbers \\
 & 27 & $+27.8$ & $-0.2$ & Gauss \\
 & 27 & $+21.6$ & $-34.5$ & Olbers \\
 & 28 & $+21.4$ & $-22.7$ & Gauss \\
 & 28 & $+0.4$ & $-11.4$ & Olbers \\
März & 2 & $-41.3$ & $+9.8$ & Gauss \\
 & 2 & $+40.1$ & $-34.9$ & Olbers \\
 & 5 & $+7.3$ & $-83.3$ & Olbers \\
 & 19 & $+56.3$ & $+36.4$ & Olbers \\
\end{tabular}
\end{center}


Bei der Rektaszension der Altonaer Beobachtung vom 2. Februar scheint ein Fehler von einer Zeitminute vorzuliegen.
